\section{Limitations and Next Steps}
The initial formulation assumes a single irreversible fraud onset, reliable
detection labels, and an observation mechanism compatible with the stated
censoring model. Platform policies may change, suspicious accounts may be
monitored more intensely, and suspensions may be noisy. The current
conditional treatment of behavioral covariates also needs an explicit audit.
These issues affect interpretation of the hazards and inferred onset times.

Before a submission is defensible, the draft needs a broader related-work
review on delayed feedback, latent event times, multi-state survival models,
and partial supervision. The proposed likelihood must be distinguished from
existing models and from the supplied conceptual note. The origin,
publication status, and appropriate attribution of that note must be
confirmed before submission. Its bibliography entry is provisional.

\paragraph{Potential impact.}
Fraud-risk estimates can generate harmful false accusations. Evaluation
should quantify false positives and calibration, examine subgroup effects
where ethically and statistically appropriate, and treat inferred onset as
uncertain model output. A research score alone should not be interpreted as
evidence sufficient for automatic punitive action.

\paragraph{Working conclusion.}
Separating occurrence from detection gives an explicit account of hidden
fraud among undetected users. Whether it yields reliable onset inference or
better early warning remains an empirical and statistical question. The
immediate goal is to determine what can be recovered under controlled delay
and supervision assumptions before expanding the architecture.
